\subsection{COMPARISON OF UNIFORMITIES}

Proposition 2 of no. 1 shows that we can take as morphisms of uniform structures the uniformly continuous mappings (Set Theory, Chapter IV, § 2, no. 1); we shall always assume in the future that the morphisms have been so chosen. In accordance with the general definitions (Set Theory, Chapter IV, § 2, no. 2), this allows us to define an order relation on the set of uniformities on a given set X:

DEFINITION 2. If \(\mathfrak{U}_1\) and \(\mathfrak{U}_2\) are two uniform structures on the same set \(\mathbf{X}\) , \(\mathfrak{U}_1\) is said to be finer than \(\mathfrak{U}_2\) (and \(\mathfrak{U}_2\) coarser than \(\mathfrak{U}_1\) ) if, denoting by \(\mathbf{X}_i\) the set \(\mathbf{X}\) with the uniform structure \(\mathfrak{U}_i\) ( \(i = 1,2\) ), the identity mapping \(\mathbf{X}_1 \to \mathbf{X}_2\) is uniformly continuous.

If \(\mathfrak{U}_1\) is finer than \(\mathfrak{U}_2\) and distinct from \(\mathfrak{U}_2\) , we say that \(\mathfrak{U}_1\) is strictly finer than \(\mathfrak{U}_2\) (and that \(\mathfrak{U}_2\) is strictly coarser than \(\mathfrak{U}_1\) ).

Two uniformities are said to be comparable if one is finer than the other.

Example. In the ordered set of uniformities on a set X, the discrete uniformity is the finest, and the coarsest uniformity is that in which the set of entourages consists of the single element \(X \times X\) .

The following proposition is an immediate consequence of Definition 1 of no. 1:

PROPOSITION 3. If \(\mathfrak{U}_1\) and \(\mathfrak{U}_2\) are two uniformities on a set \(X\) , then \(\mathfrak{U}_1\) is finer than \(\mathfrak{U}_2\) if and only if every entourage of \(\mathfrak{U}_2\) is an entourage of \(\mathfrak{U}_1\) .

COROLLARY. Let \(\mathfrak{U}_1\) and \(\mathfrak{U}_2\) be two uniformities on a set \(X\) , and suppose that \(\mathfrak{U}_1\) is finer than \(\mathfrak{U}_2\) ; then the topology induced by \(\mathfrak{U}_1\) is finer than the topology induced by \(\mathfrak{U}_2\) .

This follows immediately from the comparison of topologies in terms of neighbourhoods (Chapter I, § 2, no. 2, Proposition 3).

Remarks. 1) It can happen that a uniformity \(\mathcal{U}_1\) is strictly finer than a uniformity \(\mathcal{U}_2\) but that the two induced topologies are identical. The following example shows this:

Let \(\mathbf{X}\) be a non-empty set. For each finite partition \(\varpi = (\mathrm{A}_i)_{1\leqslant i\leqslant n}\) of \(\mathbf{X}\) , let \(\mathrm{V}_{\overline{\omega}}\) denote

\[
\bigcup_ {i} \mathrm{A} _ {i} \times \mathrm{A} _ {i}.
\]

The sets \(V_{\overline{w}}\) then form a fundamental system of entourages of a uniformity \(\mathcal{U}\) on \(X\) . For if \(\varpi\) is any finite partition of \(X\) we have \(\Delta \subset V_{\overline{w}}\) and \(V_{\overline{w}} \circ V_{\overline{g}} = \overline{V}_{\overline{w}}^{1} = V_{\overline{w}}\) (§ 1, no. 1, Example 2); and if \(\varpi' = (B_j)\) and \(\varpi'' = (C_k)\) are two finite partitions of \(X\) , then those of the sets \(B_j \cap C_k\) which are not empty form a finite partition \(\varpi\) of \(X\) , and we have \(V_{\overline{w}} \subset V_{\overline{w}'} \cap V_{\overline{w}''}\) . \(\mathcal{U}\) is called the uniformity of finite partitions on \(X\) . The topology induced by. \(\mathcal{U}\) is the discrete topology, since for each \(x \in X\) the sets \(\{x\}\) and \(C\{x\}\) form a finite partition of \(X\) . Nevertheless, if \(X\) is infinite, it is clear that \(\mathcal{U}\) is strictly coarser than the discrete uniformity. 2) If \(f: X \to X'\) is a uniformly continuous mapping, then \(f\) remains uniformly continuous if we replace the uniformity of \(X\) by a finer uniformity and that of \(X'\) by a coarser uniformity (no. 1, Proposition 2). In other words, the finer the uniformity of \(X\) and the coarser the uniformity of \(X'\) , the more uniformly continuous mappings there are of \(X\) into \(X'\) .

